\subsection{A concrete hiding and computationally binding bit commitment}\label{sec:realize:pedersen}
We use Pedersen's group commitment~\cite{Pedersen1991}
for \emph{binary coordinates}, not for raw field addition in $L$.

\begin{definition}[Prime-order commitment parameters]\label{def:realize:group}
Let $G$ be a cyclic multiplicative group of \emph{odd prime}
order $P$, with generators $g,h\in G$. Setup is trusted only
in the following precise sense: no efficient participant
knows $\alpha=\log_g h$.
The discrete-logarithm assumption is that a probabilistic
polynomial-time algorithm given the public pair $(g,h)$
cannot compute $\alpha$ except with negligible probability.
We require $P$ to exceed every parity row's number of summands.
\end{definition}

\begin{definition}[Pedersen commitment to a binary coordinate]\label{def:realize:commit}
For $b\in\{0,1\}\subset\mathbb Z_P$ and
$r\in\mathbb Z_P$, define
\[
 \mathsf{Com}(b;r):=g^b h^r\in G.
\]
An opening is the pair $(b,r)$, with $b$ restricted to
$\{0,1\}$. Every bit commitment uses a fresh independent
uniform $r$.
\end{definition}

\begin{lemma}[Exact perfect hiding]\label{lem:realize:hiding}
For each fixed bit $b$, if $r$ is uniform in $\mathbb Z_P$,
then $\mathsf{Com}(b;r)$ is uniform on $G$.
Consequently, any vector of independently blinded bit
commitments is jointly uniform on the appropriate power of $G$
and independent of its binary message vector.
\end{lemma}
\begin{proof}
Because $h$ generates $G$, the map $r\mapsto h^r$
is a bijection from $\mathbb Z_P$ to $G$.
Multiplication by the fixed $g^b$ is another bijection.
The composition sends uniform $r$ to uniform group elements.
For a vector of commitments, the maps act on independently
sampled blindings coordinatewise, so the joint distribution
is the product of the individual uniform distributions.
\end{proof}

\begin{lemma}[Two openings recover the forbidden setup logarithm]\label{lem:realize:binding-reduction}
If
$g^b h^r=g^{b'}h^{r'}$ with
$b\ne b'$ in $\mathbb Z_P$, then
$r'\ne r$ and
\[
 \boxed{\log_g h=(b-b')(r'-r)^{-1}\pmod P.}
\]
\end{lemma}
\begin{proof}
Rearrange the commitment equality:
$g^{b-b'}=h^{r'-r}=g^{\alpha(r'-r)}$.
If $r'=r$, then $g^{b-b'}=1$, implying
$b=b'$ because $g$ has order $P$.
Thus $r'-r\ne0$, and the inverse exists in
the prime field $\mathbb Z_P$. Equate the exponents modulo $P$
and divide to obtain the formula.
\end{proof}

\begin{theorem}[Computational binding for bits]\label{thm:realize:binding}
Under the discrete-logarithm assumption of
Definition~\ref{def:realize:group}, no efficient
committer can output a group element and valid
openings to two distinct bits with non-negligible probability.
\end{theorem}
\begin{proof}
Given such an algorithm, run it on
public $(G,P,g,h)$. Whenever it outputs
two distinct openings, apply
Lemma~\ref{lem:realize:binding-reduction}.
This computes $\log_g h$ in polynomial time
with the algorithm's success probability,
contradicting the discrete-logarithm assumption.
\end{proof}

\begin{definition}[Aggregating commitments]\label{def:realize:product}
For an index set $T$ of $m$ committed bits
$C_i=\mathsf{Com}(b_i;r_i)$, put
\[
 C_T:=\prod_{i\in T} C_i,\qquad
 S_T:=\sum_{i\in T}b_i\in\{0,\ldots,m\},
 \qquad R_T:=\sum_{i\in T}r_i\in\mathbb Z_P.
\]
The sum $S_T$ is an \emph{ordinary integer},
subsequently embedded in $\mathbb Z_P$.
\end{definition}

\begin{lemma}[Exact aggregate opening]\label{lem:realize:aggregate}
For each finite $T$,
\[
 \boxed{C_T=g^{S_T}h^{R_T}.}
\]
\end{lemma}
\begin{proof}
Multiply the defining equations
$C_i=g^{b_i}h^{r_i}$ over all $i\in T$.
Because $G$ is abelian, powers of $g$
and $h$ may be collected separately.
Their exponents add, yielding the identity.
\end{proof}

\begin{lemma}[Why naive XOR homomorphism fails]\label{lem:realize:no-xor}
For two committed bits $b_1=b_2=1$,
their aggregate opens to the integer $S_T=2$,
whereas $b_1+b_2=0$ in $\F_2$.
Consequently the natural product of two
bit commitments is not automatically an
opening of their XOR with the sum of blindings.
\end{lemma}
\begin{proof}
The product is $g^2h^{r_1+r_2}$
by Lemma~\ref{lem:realize:aggregate}.
An opening to the XOR using the same
blinding $r_1+r_2$ would instead be
$h^{r_1+r_2}$. As $g$ has odd prime order
$P$, we have $g^2\ne1$. Thus the elements
are different.
\end{proof}

\begin{remark}[Setup and post-quantum scope]
The Pedersen instantiation in a conventional
discrete-logarithm group is \emph{not}
post-quantum secure against general quantum
computers. It is an explicit classical
commitment/proof candidate for the affine
checker, not a post-quantum implementation
of the original KBFold PCS. A transparent
post-quantum realization remains separate.
\end{remark}
